\section{Bootstrapping basic constructs}
\label{bootstrapping-basic}
The base syntax of ILD is insufficient for writing nontrivial programs
in a succinct way. However, its \(\ildsf{macroexpand}\) facility allows
extending the syntax with arbitrary metaprogramming constructs, which we
build incrementally.

\begin{note}
maybe here add a note that this power is similar to fexpr stuff but uses a different mechanism?
\end{note}

\subsection{Lambda}
\label{lambda-macro}
First, we define a macro called \ild{lambda} that will let us introduce new
abstraction objects (\cref{abstraction}). We want expansion of the \ild{lambda}
macro to result in a call to \ild{mk-lambda} given the current binding environment
(\ild{free-vars}) and the quoted argument names and body:
\begin{ildcode}
(!lambda (x y) (+ x y))
; should become
(mk-lambda (free-vars) (quote (x y))
  (quote (+ x y)))
\end{ildcode}
The outer call to \ild{mk-lambda} introduces a new \(\ildsetAbstr\)-object
with the current binding environment%
\footnote{The \ild{(free-vars)} closure given to the host \ild{mk-lambda} can be replaced
by a closure that contains only \ild{cons}, \ild{mk-lambda}, \ild{quote} and \ild{free-vars},
since those are the only bindings used by the body of the macro.}
and a body that embeds the \emph{value} of \ild{mk-lambda}, to in turn
expand to a \ild{mk-lambda}-expression when called:
\begin{ildcode}
(mk-lambda
  (free-vars)
  '(arg-names body)
  '(cons mk-lambda
    (cons (cons free-vars '())
      (cons (cons quote (cons arg-names '()))
        (cons (cons quote (cons body '()))
          '())))))
\end{ildcode}
\lstcaption{Implementation of the lambda macro}
Embedding the values of \ild{mk-lambda}, \ild{free-vars} and \ild{quote}
avoids name shadowing problems traditionally encountered with unhygienic
macro systems (\cref{values-in-macroexpand}).

This definition can then be wrapped in another call of \ild{mk-lambda} in order to expose
it under a \ild{lambda} name usable from the scoped body -- supposing that
this definition is stored as an ILD program called \texttt{core/bootstrap/lambda-macro.ild},
we recall it twice (once to define \ild{lambda}, and once to pass itself as the
value of \ild{lambda}):

\begin{ildcode}
((!(eval (free-vars)
    (read-source
      "core/bootstrap/lambda-macro.ild"))
 (lambda)

  ...) ; body that may use the lambda macro

  ; definition of lambda
  (eval (free-vars)
    (read-source
      "core/bootstrap/lambda-macro.ild")))
\end{ildcode}
\lstcaption{Outermost usage of the lambda macro}

\subsection{Basic utilities}
We define helper utilities in a scope where the \ild{lambda} macro is present.
To keep the code succinct, we store these definitions in separate bootstrap
source files and introduce them in scope with the same technique as the
\ild{lambda} macro.
\begin{ildcode}
; definition of list
(!lambda args args)

; definition of cadr
(!lambda (p) (car (cdr p)))

; Y is special case of poly-fix
(!lambda (f) (car (poly-fix f)))
\end{ildcode}
\lstcaption{Definitions of basic helper primitives}
To more easily write macros that generate expressions
of the form \ild{(!lambda <args> <body>)}, we define a helper:
\begin{ildcode}
(!lambda (args body)
  (cons macroexpand
    (cons lambda
      (cons args
        (cons body '())))))
\end{ildcode}
\lstcaption{Definition of \ild{expand-lambda}}
Branching in ILD is done via the \ildhost{bool-to-k} host
function. In order to only evaluate the selected branch, we delay
evaluation to after the condition is known by wrapping in a 0-argument
function, and force it afterwards:
\begin{ildcode}
(!lambda (cond t f)
  ; generate an expression that
  (list     ; forces the result
    (list   ; applies the selector,
            ; obtaining a delayed value
      (list bool-to-k cond) ; selector

      ; delayed #t and #f branches
      (expand-lambda '() t)
      (expand-lambda '() f))))
\end{ildcode}
\lstcaption{Definition of the \ild{if} macro}

\begin{ildcode}
(Y (!lambda (map)
  (!lambda (f l)
    (!if (null? l)
      l
      (!if (pair? l)
        (cons
          (f (car l)) (map f (cdr l)))
        (make-fail
          (list 'not-a-list l)))))))
\end{ildcode}
\lstcaption{Definition of the \ild{map} function}

\subsubsection{Let}
The binding construct \ild{let} is simply a more ergonomic way
of writing \ild{lambda} with specifically-named arguments:
\begin{ildcode}
(!lambda (letlist body)
  (cons
    (expand-lambda
      (map car letlist)
      body)
    (map cadr letlist)))
\end{ildcode}
\lstcaption{Definition of \ild{let}}
Now we can bind locally-scoped values ergonomically:
\begin{ildcode}
((!lambda (add1 fourty-two)
    (add1 fourty-two))

    (!lambda (x)
      (add x 1)) ; definition of add1
    42)          ; definition of fourty-two
; produces 43
\end{ildcode}
becomes
\begin{ildcode}
(!let (
    (add1 (!lambda (x)
            (add x 1)))
    (fourty-two  42))

  (add1 fourty-two))
\end{ildcode}
\lstcaption{Rewriting a \ild{lambda}-based scope binding to use \ild{let}}

\subsection{Letrec}
\label{letrec}
We define letrec as a macro that builds each binding as a recursive operator
that in turn takes all bindings as arguments, and passes the list of
those bindings to \ild{poly-fix}:

\begin{ildcode}
(!lambda (defs body)
  (!let
    ((args (map car defs))
      (item-bodies (map cadr defs)))
    (!let (
      (arg-bodies
        (map
          (!lambda (def-body)
            (expand-lambda args def-body))
          item-bodies))
      (flist (gensym "flist")))

      (!let (
        (result-body (cons
          (expand-lambda args body)
          (generate-element-getters
            flist args))))

        (list
          (expand-lambda
            (list flist)
            result-body)
          (cons poly-fix arg-bodies))))))
\end{ildcode}
\lstcaption{Definition of letrec}

We omit the definition of the \ild{generate-element-getters} helper, which is merely technical.
